% ============================================================
%  Mehdi Eskandari-Ghadi — Resume
%  Compile: pdflatex resume.tex  (or use the Makefile)
% ============================================================
\documentclass[10pt, letterpaper]{article}

% --- Packages -----------------------------------------------
\usepackage[margin=0.65in, top=0.55in, bottom=0.55in]{geometry}
\usepackage{enumitem}
\usepackage{tabularx}
\usepackage{array}
\usepackage{fontenc}
\usepackage{inputenc}
\usepackage[T1]{fontenc}
\usepackage{mathptmx}        % Times-like serif (similar to Cambria)
\usepackage{hyperref}
\usepackage{xcolor}
\usepackage{parskip}

% --- Hyperlink styling --------------------------------------
\hypersetup{
  colorlinks=true,
  urlcolor=black,
  linkcolor=black
}

% --- Layout helpers -----------------------------------------
\setlength{\parindent}{0pt}
\setlength{\parskip}{0pt}

% Section header with full-width rule beneath
\newcommand{\rsection}[1]{%
  \vspace{8pt}%
  {\large\bfseries\MakeUppercase{#1}}%
  \vspace{1pt}%
  \hrule height 0.8pt%
  \vspace{4pt}%
}

% Job/role line: title left, date right
\newcommand{\jobline}[2]{%
  \vspace{5pt}%
  \noindent\textbf{#1}\hfill\textbf{#2}\par%
  \vspace{1pt}%
}

% Subline (institution, location — plain italic)
\newcommand{\subline}[1]{%
  \vspace{-1pt}%
  \noindent\textit{#1}\par%
  \vspace{1pt}%
}

% Bullet list settings
\newenvironment{rlist}{%
  \begin{itemize}[leftmargin=1.4em, itemsep=1pt, topsep=2pt, parsep=0pt]%
}{%
  \end{itemize}%
}

% ============================================================
\begin{document}
\pagestyle{empty}

% --- Header -------------------------------------------------
\begin{center}
  {\LARGE\bfseries Mehdi Eskandari-Ghadi}\\[4pt]
  mehdi.esk73@gmail.com\ $|$\ (303)378-1407\ $|$\ Salt Lake City, UT\\[2pt]
  \href{http://www.linkedin.com/in/M-E-Ghadi}{www.linkedin.com/in/M-E-Ghadi}\ $|$\
  Authorized to work in the U.S.\ $|$\ Sponsor not required
\end{center}

\vspace{2pt}

% --- Education ----------------------------------------------
\rsection{Education}

\jobline{Civil Engineering Ph.D. (GPA 4.0)}{Boulder, CO}
Engineering Science Focus \hfill 2016 -- 2021\\
University of Colorado Boulder

% --- Technical Skills ---------------------------------------
\rsection{Technical Skills}

\vspace{2pt}
\begin{tabular}{@{}p{4cm}p{13cm}@{}}
  \textbf{Programming}       & Python, TypeScript/JavaScript, VBA, MATLAB, AutoLISP \\[2pt]
  \textbf{Platforms \& Data} & Palantir Foundry, Apache Spark, Polars \\[2pt]
  \textbf{Engineering}       & AutoCAD, PLSCADD, L-PILE, ABAQUS/LS-DYNA \\[2pt]
  \textbf{Standards \& Tools} & NESC, GO-95, FERC-881, Git, Microsoft Office Suite, LaTeX \\
\end{tabular}

% --- Engineering Experience ---------------------------------
\rsection{Engineering Experience}


\jobline{Reliability \& Standards Software Engineer II}{Jan 2025 -- Present}
\subline{PacifiCorp}
\begin{rlist}
  \item Built full-stack engineering workflows on the Palantir Foundry platform using TypeScript and Python, covering audit preparation, data management UIs, interdepartmental data collection, and automated queue-based processing pipelines.
  \item Designed a git-inspired facility lifecycle management system within Foundry for branching, modifying, and publishing transmission facilities, with automated resource tracking, branch conflict resolution, and staged approval workflows.
  \item Architected backend data pipelines in Python using Spark and Polars for large-scale ingestion, parsing, and health remediation of legacy spreadsheets, automating correction workflows and multi-source data merging to bring records to quality standards.
  \item Implemented IEEE-compliant thermal rating models in Python covering solar heat gain, convective cooling, radiative heat loss, and conductor ampacity-resistance energy balance, enabling dynamic and static line rating computations within the platform.
  \item Implemented Python and TypeScript routines within Foundry to perform geometric analysis of parallel transmission corridors, computing mutual impedance candidate pairs and characterizing average and percentile conductor separation distances, with line-to-line distance profiles visualized in platform-integrated UI graphs.
  \item Overhauled legacy spreadsheets with VBA-based functionalities, maintainable assumption/reference sheets, and improved UX, establishing a more auditable and automatable calculation baseline ahead of platform migration.
  \item Developed AutoLISP routines for AutoCAD one-line drawings that parse shape geometries and layer data to identify transmission lines, facilities, and equipment; deduce line segments and breaker-to-breaker sections including different equipment configurations; and export structured data to spreadsheets to bring drawing databases and engineering records to convergence.
\end{rlist}

\jobline{Transmission Engineer I/II}{Jan 2023 -- Jan 2025}
\subline{PacifiCorp}
\begin{rlist}
  \item Utilizing high-level knowledge of transmission line design to scope new transmission lines in collaboration with other engineering disciplines.
  \item Utilizing in-depth knowledge of PLSCADD, L-PILE, NESC, and GO-95 for design of new power transmission lines and rebuilding/maintenance of existing lines in states of UT, WY, OR, WA, and CA in collaboration with teams of engineers, managers, field crew, and legal.
  \item Utilizing in-depth knowledge of structural engineering and geotechnical engineering for custom-engineered solutions of transmission structures and drilled pier foundations with reference to ASCE, ACI, and FHWA to meet special needs.
  \item Reduced engineering hours by threefold and cost by approx.\ \$1.5M by intersecting programming and automation experience and transmission engineering experience to create a line rating analysis tool. This analysis tool allows compliance with FERC-881 with minimum reduction of line capacity, increased calculation accuracy, and reduced engineering time.
  \item Utilized visual basic programming skills to create automated tools for data analysis and modeling of existing transmission line ratings to assess FERC-881 compliance.
  \item Applying programming and software design skills and IEEE line rating knowledge to guide the development of proprietary software tools for FAC-008 compliance.
  \item Employed excel programming and data analysis skills to develop automated tools for deep foundation design, construction package generation, and transmission pole design aids.
  \item Utilized transmission engineering knowledge and programming experience to generate company sag-tension standards for tree-wire.
  \item Designing and reviewing drawings for tubular steel poles and drilled piers to withstand vertical and lateral loads.
  \item Performing detail-oriented technical review of transmission standards and collaboration in their developments.
  \item Initiating process improvement pipelines for transmission design processes and documentations.
  \item Improved company standards and company procedures with detail-oriented reviews/QAQC.
\end{rlist}

% --- Research Experience ------------------------------------
\rsection{Research Experience}

\noindent\textbf{Academic Research}

\jobline{Graduate Research Assistant / Ph.D. Candidate}{Jun 2018 -- Jan 2023}
\subline{University of Colorado Boulder, Boulder, CO}
\begin{rlist}
  \item Collaborated with a team of three senior researchers at the Lawrence Berkeley National Laboratory to investigate and develop mechanistic models for subcritical crack growth healing.
  \item Utilized knowledge of fundamentals of fracture mechanics, fluid transport, and boundary element methods to develop a software that couples fields of mechanics and physics.
  \item Merged expertise in research, programming, and engineering to investigate multiphysical and multiscale solid-fluid interactions (e.g., adsorption, surface tension, solvation pressure) in porous solids and develop a thermodynamics-poromechanics-based MATLAB software for ``adsorption-induced deformation of microporous material''.
  \item Employed technical writing and illustrative skills for preparation of 4 peer-reviewed publications in journals of JMPS, IJSS, and EFM.
\end{rlist}

\vspace{4pt}
\noindent\textbf{Talks and Presentations}

\jobline{EMI Virtual Conference 2021 and EMI Conference 2022}{May 2021 and May 2022}
\begin{rlist}
  \item Applied illustrative and verbal communication skills to deliver presentations titled ``A multi-scale theory explaining the initial shrinkage of micro-porous solid upon gas adsorption'' and ``A Multiphysical Surface-Force Based Fracture Theory for Subcritical Crack Growth in Surface-Reactive Environments'', followed by Q\&A from engineering and academic audience.
  \item Received runner up award in the 2021 Student Poster Competition for poster presentation on sorption-induced deformation of porous material.
\end{rlist}

\jobline{Lawrence Berkeley National Laboratory BES}{Sep 2020 -- Jan 2023}
\begin{rlist}
  \item Presented regular scientific presentations on topic of crack propagation to teams of research professionals and senior scientists in the LBNL BES with diverse backgrounds in mechanics, physics, and chemistry.
\end{rlist}

\jobline{Extracurricular Projects}{2018}
\begin{rlist}
  \item Programmed a discrete element MATLAB code for modeling of grain structure stability and avalanche.
\end{rlist}

% --- Teaching Experience ------------------------------------
\rsection{Teaching Experience}

\noindent\textbf{Teaching Assistantships/Tutoring}

\jobline{Teaching Assistant / Tutor}{Aug 2016 -- Dec 2019}
\subline{University of Colorado Boulder, Boulder, CO}
\begin{rlist}
  \item Presented class lectures and laboratory teaching sessions, held hours of office hours and exam preparation sessions, prepared and graded exams and homework question sets, for courses of Advanced mechanics of materials (graduate level) and Dynamics (undergraduate level).
  \item Tutored students for FE exam with 100\% success rate.
  \item Tailored technical communication specific to each student's understanding and utilized CAD skills to enhance visualization and understanding.
\end{rlist}

\jobline{Teaching Assistant}{Fall 2014 and Fall 2015}
\subline{University of Tehran, Tehran, Iran}
\begin{rlist}
  \item Held weekly office hours and graded homework, quizzes, and class projects, for courses of Structural Analysis II (undergraduate level) and Hydraulics of open channels (undergraduate level).
\end{rlist}

% --- Leadership Roles ---------------------------------------
\rsection{Leadership Roles}

\noindent\textbf{Mentorship}

\jobline{Research Mentor}{Jan 2022 -- Jan 2023}
\subline{University of Colorado Boulder, Boulder, CO}
\begin{rlist}
  \item Trained new Ph.D.\ students on fundamentals of mechanics, written and oral skills, and technical research.
  \item Guided a research trainee through the modeling and simulation of beam bending in humidity actuators.
\end{rlist}

\jobline{President/Co-president --- Persian Student Organization, Boulder, CO}{Aug 2018 -- May 2019}
\subline{University of Colorado Boulder, Boulder, CO}
\begin{rlist}
  \item Proposed, secured funding, planned, and oversaw student events in collaboration with professionals from industry and academia, local businesses, and non-profit organizations.
\end{rlist}

% --- Other Interests ----------------------------------------
\rsection{Other Interests}

\begin{rlist}
  \item I am a plant hobbyist, animal lover, dog owner, and skilled chef for humans and dogs.
\end{rlist}

\end{document}